\documentclass[10pt,a4paper]{amsart}
\usepackage[margin=19mm]{geometry}
\usepackage{amsmath,amssymb,mathtools}
\usepackage[scr=rsfs,cal=euler]{mathalfa}
\usepackage{xfrac}
\usepackage[unicode,hidelinks,bookmarks=false]{hyperref}
\newcommand{\ie}{i.e.\ }
\newcommand{\cf}{cf.\ }
\newcommand{\resp}{respectively\ }
\newcommand{\Subsection}{\subsection}
\providecommand{\nobreakdash}{\nobreak\hbox{-}}
\setlength{\emergencystretch}{2em}
\multlinegap0pt
\usepackage{bm}
\usepackage{bbm}
\usepackage{dsfont}
\usepackage{xspace}
\usepackage{cancel}
\usepackage{mathtools}
\usepackage{amsthm}
\makeatletter
\@ifundefined{lemma}{\newtheorem{lemma}{Lemma}}{}
\@ifundefined{theorem}{\newtheorem{theorem}[lemma]{Theorem}}{}
\makeatother
\newcommand{\rC}{\mathrm{C}}
\title[Rigidity of operator systems: tight extensions and noncommut]{Rigidity of operator systems: tight extensions and noncommutative measurable structures}
\author{Rapha\"el Clou\^atre, Ian Thompson}
\date{Source: arXiv:2406.16806v2 (CC BY 4.0)}
\begin{document}
\maketitle

\noindent\emph{Source and licence.} Rigidity of operator systems: tight extensions and noncommutative measurable structures. Version arXiv:2406.16806v2 (CC BY 4.0), distributed under the Creative Commons Attribution 4.0 International licence, as recorded in the arXiv licence metadata for this exact version and shown on the abstract page https://arxiv.org/abs/2406.16806v2. Full rights notice: licenses/arxiv-2406.16806-CC-BY-4.0.txt.\par\smallskip

\noindent\textit{Editorial scope.} Counted unit: the Theorem whose source label is \texttt{T:normhomo} in the article (source0.tex lines 818--820, source label \texttt{T:normhomo}), delivered together with the article's own complete original proof of it (source0.tex lines 821--870, one \texttt{proof} environment of 50 lines). No mathematics is written, repaired or re-derived: statement and proof are the article's own text line for line. Required context retained uncounted: none. Every definition, display and label of those lines is retained unchanged. Omitted from this package and not redistributed: 1--817, 871--954. Nothing inside a retained excerpt is omitted or altered: every retained body is delivered exactly as the article's lines are written, so no omission receipt of any kind accompanies this package. Citations: the retained excerpts carry no citation command at all, so no bibliography entry of the article is transcribed and no reference is redistributed. Reference adaptation: none was needed and none was made; every reference inside the delivered unit resolves inside it, so no unresolved reference or citation is produced. Printed slips kept literal: no printed word, symbol, hypothesis or number of the counted statement or of its proof was altered. Layout and class: the article's own class is \texttt{amsart} with packages including graphicx, amscd, amsmath, amsthm, amsfonts, amssymb, mathrsfs, bm, enumerate, amsrefs, xcolor, hyperref, inputenc, tikz; the wrapper is the lane's pinned \texttt{amsart} editorial wrapper at 10pt on A4, which restates the theorem-like environments the retained text needs and adds the article's own macro definitions that the retained text needs (\texttt{\textbackslash rC}), with extra wrapper packages (bm, bbm, dsfont, xspace, cancel, mathtools). The delivered numbering is the unit's own. Assets: no retained span contains a figure environment, a graphics-inclusion command, a PDF-page inclusion, a table or a TikZ picture; no image file is copied into this package, the package declares no asset dependency and no image file is redistributed with it. Licence: version arXiv:2406.16806v2 of this article is distributed under the Creative Commons Attribution 4.0 International licence, as recorded in the arXiv licence metadata for this exact version and shown on the abstract page https://arxiv.org/abs/2406.16806v2. A full rights notice is delivered at licenses/arxiv-2406.16806-CC-BY-4.0.txt. Characters: every retained excerpt is pure ASCII, so the delivered engine, XeLaTeX with its default Latin Modern text font, reports no missing character.
\begin{theorem}\label{T:normhomo}
Let $A$ be a unital  $\rC^*$-algebra. Let $\pi:A\to B(H)$ be a unital $*$-representation and let $\Upsilon_\pi$ denote the universal representation of $\pi(A)$. If $\pi(A)$ is homogeneous and $\Upsilon_\pi\circ \pi$ has the unique tight extension property, then $\pi$ has the strong Korovkin rigidity property with respect to $S$.
\end{theorem}

\begin{proof}
Since $\pi(A)$ is homogeneous, there is a positive integer $d$ such that every irreducible $*$-representation of $\pi(A)$ is $d$-dimensional. Fix once and for all a Hilbert space $F$ of dimension $d$. Then, every irreducible $*$-representation of $\pi(A)$ is unitarily equivalent to an irreducible $*$-representation on $F$. 

Let $\psi_\gamma:A\to \pi(A)$ be a net of unital completely positive maps such that 
\[
\lim_\gamma\|\psi_\gamma(s)-\pi(s)\|=0, \quad s\in S.
\]
Fix a self-adjoint element $a\in A$. The goal is to show that the net $(\psi_\gamma(a))$ converges to $\pi(a)$ in norm. Equivalently, we must show that any subnet of $(\psi_\gamma(a))$ has a further subnet converging in norm to $\pi(a)$.

Fix  a subnet  $(\psi_\lambda(a))$ of $(\psi_\gamma(a))$. For each $\lambda$, by the first paragraph of the proof we may find an irreducible $*$-representation $\sigma_\lambda:\pi(A)\to B(F)$ such that 
\begin{equation}\label{Eq:normrep}
\|\sigma_\lambda(\psi_\lambda(a)-\pi(a))\|=\|\psi_\lambda(a)-\pi(a)\|.
\end{equation}
Let $\sigma:\pi(A)\to B(F)$ be a cluster point of the net $(\sigma_\lambda)$ in the pointwise norm topology, so that there is a subnet $(\sigma_{\lambda_\mu})$ satisfying
\begin{equation}\label{Eq:conv1}
\lim_\mu\|\sigma(\pi(b))-\sigma_{\lambda_\mu}(\pi(b))\|=0, \quad b\in A.
\end{equation}
Then, $\sigma$ is a unital $*$-representation of $\pi(A)$, and it must necessarily be irreducible since $\pi(A)$ is homogeneous.



We claim that 
\begin{equation}\label{Eq:conv2}
\lim_\mu \|\sigma_{\lambda_\mu}(\psi_{\lambda_\mu}(b))-\sigma(\pi(b))\|=0, \quad b\in A.
\end{equation}
To see this, consider the weak-$*$ continuous $*$-representation $\widehat\sigma:\pi(A)^{**}\to B(F)$ and let $z\in \pi(A)^{**}$ be the central projection such that $\ker \widehat\sigma=\pi(A)^{**}(1-z)$. Since $F$ is finite-dimensional and $\sigma$ is irreducible, we have $B(F)=\sigma(\pi(A))$ and there is a weak-$*$ homeomorphic $*$-isomorphism $\Omega:B(F)\to \pi(A)^{**}z$ such that $\Omega(\sigma(\pi(a)))=\pi(a)z$ for each $a\in A$.  For each $\mu$, we define a unital completely positive map $\Psi_\mu:A\to \pi(A)^{**}$ as
\[
\Psi_\mu(b)=(\Omega\circ \sigma_{\lambda_\mu} \circ \psi_{\lambda_\mu})(b)\oplus \pi(b)(1-z)
\]
for every $b\in A$. By construction, for each $s\in S$ we see that
\[
\lim_\mu \|\sigma_{\lambda_\mu}(\psi_{\lambda_\mu}(s))-\sigma(\pi(s))\|=0
\]
and so
\[
\lim_\mu \|\Psi_\mu(s)-\pi(s)\|=0.
\]
Let $\Phi:A\to \Upsilon_\pi(A)''$ be a cluster point of the net $(\Upsilon_\pi\circ \Psi_\mu)$ in the pointwise weak-$*$ topology. Then, $\Phi$ is a unital completely positive map agreeing with $\Upsilon_\pi\circ \pi$ on $S$. By the unique tight extension property of $\Upsilon_\pi\circ \pi$, we infer that $\Phi=\Upsilon_\pi\circ \pi$ on $A$. Applying $\widehat{\Upsilon_\pi}^{-1}$, we find that the net $(\Psi_\mu(b))$ converges to $\pi(b)$ in the weak-$*$ topology of $\pi(A)^{**}$ for every $b\in A$. In particular,  for $b\in A$ the net $((\Omega\circ \sigma_{\lambda_\mu} \circ \psi_{\lambda_\mu})(b))$ converges to $\pi(b)z$  in the weak-$*$ topology of $\pi(A)^{**}$ and, in turn, by applying $\Omega^{-1}$ we see that $((\sigma_{\lambda_\mu} \circ \psi_{\lambda_\mu})(b))$ converges to $\sigma(\pi(b))$  in the norm topology of $B(F)$ (since $F$ is finite-dimensional). This establishes the claim.

Finally, upon combining \eqref{Eq:conv1} and \eqref{Eq:conv2}, we find 
\[
\lim_\mu \|\sigma_{\lambda_\mu}(\pi(a)-\psi_{\lambda_\mu}(a))\|=0.
\]
By \eqref{Eq:normrep}, this implies 
\[
\lim_\mu \|\pi(a)-\psi_{\lambda_\mu}(a)\|=0
\]
as desired.

\end{proof}

\end{document}
